\subsection*{把长表达式拆成三个短问题}
考虑 $z=wa+b$、输出 $\hat y=\sigma(z)$、损失 $L=(\hat y-y)^2$。若直接对嵌套式求导，容易漏项。先分别问：损失对输出多敏感，输出对预激活多敏感，预激活对权重多敏感？答案是 $2(\hat y-y)$、$\sigma'(z)$、$a$。

把三个响应倍率按路径相乘，就得到权重梯度。输入激活 $a$ 越大，同一权重变化越能改变预激活；激活函数若进入饱和区，$\sigma'(z)$ 很小，梯度会被削弱。公式因此也解释了训练现象，而不仅给出可执行运算。

\subsection*{先检查最简单的线性输出}
让 $\sigma$ 为恒等函数，取 $a=2,b=1,y=4,w=1$。预测为三，误差为负一，权重梯度为负四。梯度下降减去负数，使权重增加，预测向目标移动，符号与直觉一致。偏置梯度为负二，因为偏置改变一个单位只使预测改变一个单位。

对多层情形，上一层的一个单元通常影响下一层多个单元，因而误差信号需要求和；矩阵转置正好把下一层的影响汇集回上一层。权重梯度的外积则把“接收端敏感性”和“发送端激活”配成所有连接的贡献。

推导时先固定向量是行还是列，再检查每个矩阵形状。损失是否带二分之一、是否按样本平均，也会改变系数。用一个可手算的标量例子核对符号和尺度，往往比只盯着大矩阵更有效。
